\chapter{System Design and Implementation}

The system was designed around one principle: avoid invoking the LLM when the
text layer alone is sufficient. This chapter documents how that principle
shapes each component, from the page classifier that decides when the LLM is
needed to the postprocessing pipeline that cleans up what it produces.

\section{Architecture Overview}

The system takes a PDF file as input and produces one combined Markdown document as output.
It consists of four modules: a page content classifier, a PDF processing module, an LLM integration layer, and a post-processing pipeline.
Each module is implemented as a separate Python package and communicates through well-defined data structures.
This keeps the components interchangeable and makes them testable without running the full pipeline.

Processing begins when a PDF is loaded by PyMuPDF \autocite{pymupdf}.
For each page, the classifier analyses structural signals and assigns one of six content types: TEXT, IMAGE, FORMULA, TABLE, MIXED, or EMPTY.
This classification determines which extraction path the page takes.
The decision is made per page, so a single document can route different pages to different strategies.

TEXT pages are handled entirely by pymupdf4llm \autocite{pymupdf4llm}, which extracts the text layer and converts it to Markdown without invoking the LLM.
All other page types are rendered as a base64-encoded PNG (Portable Network Graphics) image and passed to a vision-capable LLM.
The LLM receives a structured prompt suited to the detected content type and returns a Markdown string.
EMPTY pages are skipped, returning a short placeholder without any processing.

Pages are converted in parallel using a thread pool. The default concurrency is four workers, which overlaps the I/O wait time of LLM calls without overloading the LM Studio server. Each worker opens its own \texttt{fitz.Document} handle, because PyMuPDF document objects are not thread-safe. The results are collected in page order and joined into a single Markdown document, with a horizontal rule between each page.

The LLM integration layer communicates with LM Studio \autocite{lmstudio} over a local HyperText Transfer Protocol (HTTP) server that exposes an OpenAI-compatible API \autocite{khan2024}. All LLM calls are routed through a single \texttt{call\_llm} function, which accepts the server URL, model name, message list, temperature, and token limit as parameters. Every strategy accepts this function as an injectable parameter, so tests can replace it with a fixed-response stub without touching the rest of the code.

The post-processing pipeline receives the joined Markdown and applies a fixed sequence of cleanup passes before writing the final output file.

Figure~\ref{fig:architecture} shows the complete processing pipeline.

\begin{figure}[H]
    \centering
    % System architecture — mirrors run() in app.py (adaptive pipeline)
\begin{tikzpicture}[
    font=\sffamily\small,
    >={Latex[length=1.8mm]},
    box/.style={draw=black!70, line width=0.5pt, rounded corners=2pt, align=center, inner sep=4pt},
    startstop/.style={box, fill=black!12, rounded corners=5pt, font=\sffamily\footnotesize\bfseries},
    proc/.style={box, fill=white, font=\sffamily\footnotesize},
    cond/.style={box, fill=black!5, font=\sffamily\footnotesize},
    note/.style={font=\sffamily\scriptsize, text=black!60},
    lbl/.style={font=\sffamily\scriptsize\itshape, inner sep=2pt}
]
\node[startstop] (pdf) at (5.5,0) {PDF file};

\node[proc, text width=42mm] (lang) at (2.4,-1.55) {language detection\\{\scriptsize\ttfamily langdetect}{\scriptsize, first page}};
\node[proc, text width=50mm] (bulk) at (8.6,-1.55) {bulk text extraction\\{\scriptsize\ttfamily pymupdf4llm}{\scriptsize, once per document}};

\node[proc, text width=62mm] (pool) at (5.5,-3.15) {thread pool -- 4 workers by default\\{\scriptsize results collected in page order}};

% --- per-page processing frame ---
\draw[black!50, dashed, thin, rounded corners=3pt] (-0.9,-4.1) rectangle (12.4,-9.55);
\node[note, anchor=north west] at (-0.9,-4.12) {per page, in parallel};

\node[cond, text width=44mm] (toc) at (5.5,-5.0) {TOC page with PDF bookmarks?};
\node[proc, text width=34mm] (tocout) at (1.1,-6.55) {TOC built from\\\texttt{doc.get\_toc()} -- no LLM};
\node[proc, text width=44mm] (clf) at (5.5,-6.55) {page classifier + routing\\{\scriptsize Figures~\ref{fig:classifier} and~\ref{fig:routing}}};

\node[proc, text width=30mm] (tp) at (1.6,-8.55) {text path:\\pre-extracted Markdown};
\node[proc, text width=40mm] (vp) at (5.5,-8.55) {vision path: 150\,DPI PNG\\\(\rightarrow\) LM Studio LLM};
\node[proc, text width=26mm] (pl) at (9.9,-8.55) {EMPTY:\\placeholder};

\draw[->] (pdf) -- (lang);
\draw[->] (pdf) -- (bulk);
\draw[->] (lang) -- (pool);
\draw[->] (bulk) -- (pool);
\draw[->] (pool) -- (toc);
\draw[->] (toc) -- node[lbl, above] {yes} (tocout);
\draw[->] (toc) -- node[lbl, right] {no} (clf);
\draw[->] (clf) -- (tp);
\draw[->] (clf) -- (vp);
\draw[->] (clf) -- (pl);

% --- document assembly ---
\node[proc, text width=46mm] (clean) at (5.5,-10.55) {\texttt{clean\_page} -- per-page cleanup};
\node[proc, text width=52mm] (join) at (5.5,-11.9) {join pages with \texttt{-{}-{}-} separator};
\node[proc, text width=58mm] (post) at (5.5,-13.25) {\texttt{postprocess\_markdown} -- document-wide cleanup};
\node[startstop] (out) at (5.5,-14.7) {Markdown output};

\draw[->] (tp) -- (clean);
\draw[->] (vp) -- (clean);
\draw[->] (pl) -- (clean);
\draw[->] (clean) -- (join);
% TOC pages bypass clean_page and go straight into the join
\draw[->] (tocout.south) .. controls (-1.6,-10.9) and (-0.5,-11.9) .. (join.west);
\draw[->] (join) -- (post);
\draw[->] (post) -- (out);
\end{tikzpicture}

    \caption[System architecture: from PDF input to Markdown output]{System architecture: from PDF input to Markdown output. TOC pages are rebuilt from the PDF bookmark outline and skip both the LLM and the per-page cleanup.}
    \label{fig:architecture}
\end{figure}

\section{Page Content Classifier}

The classifier extracts four structural signals from each page using PyMuPDF \autocite{pymupdf}: the number of embedded images, the number of vector drawing paths, the total character count of the text layer, and the number of detected tables.
No machine learning is used. The classification is entirely rule-based, which keeps it fast and predictable.

Formula detection uses two independent checks.
First, the classifier counts vector paths whose bounding box is smaller than 20 pixels in either dimension.
More than 30 such short paths on a single page indicates the presence of rendered mathematical symbols, since PDF renders formula components as small vector segments.
Second, the classifier scans the text layer for mathematical Unicode characters such as $\sum$, $\int$, and $\nabla$, as well as LaTeX-style sequences that OCR engines sometimes preserve.
Either check alone is sufficient to flag the page as formula-containing.

The classification follows a fixed priority order, working from the narrowest
conditions toward the broader ones. The first check is for empty pages: fewer
than ten characters, no images, and fewer than five vector paths. The next
checks target the most visually distinctive content. A formula page has formula
signals but fewer than 50 characters, since formula-only layouts tend to carry
little surrounding body text. An image page has more than three embedded images,
and a table page has at least one table detected by PyMuPDF's table finder.
Pages that carry image or formula signals but did not qualify for any of those
specific types fall into MIXED. Plain-text pages, those with 50 or more
characters and no visual signals, are classified as TEXT. Any page that passes
every check without a clear match falls back to MIXED with a confidence score
of 0.5.

Each classification carries a confidence score between 0.5 and 0.95.
TEXT pages score the highest at 0.95. To reach that classification, a page must have at least 50 characters and no image, formula, or table signals. That is the hardest combination to satisfy by accident.
EMPTY and IMAGE pages score 0.9. EMPTY requires fewer than ten characters, no images, and fewer than five vector paths, which is a very narrow condition. IMAGE requires more than three embedded images, which rarely appears on a text-heavy page.
FORMULA and TABLE score 0.85 because both detectors can misfire. Vector paths above the threshold sometimes come from decorative page elements rather than formulas. PyMuPDF's table finder occasionally flags dense grid layouts that are not real tables.
MIXED pages score either 0.75 or 0.5. Pages that carry image or formula signals but did not qualify for IMAGE, FORMULA, or TABLE --- for example, a page with one or two embedded images alongside substantial text --- score 0.75. The fallback MIXED at 0.5 catches the remaining pages: fewer than 50 characters with no images, no formulas, and no detected tables. No signal was strong enough to assign a more specific type.
The confidence score is stored in the \texttt{PageAnalysis} result but plays no role in the current routing logic. Pages are routed by type label alone, regardless of score. This is intentional: the rule-based classifier is deterministic, so the same signals always produce the same type and a score threshold would not change any routing decision. A future version could use the score to fall back to the hybrid strategy on genuinely ambiguous pages, where the MIXED 0.5 fallback fires.

Figure~\ref{fig:classifier} shows the classification cascade. Figure~\ref{fig:routing} shows the routing checks that follow it.

\section{PDF Processing Module}

The PDF processing module has two sub-pipelines: one that extracts the text layer and one that renders pages as images.
Which pipeline runs depends on the page type assigned by the classifier.
TEXT pages go through text extraction only. IMAGE, FORMULA, and MIXED pages go through image rendering. EMPTY pages skip both.
The hybrid strategy uses both pipelines on the same page.

TABLE pages are a special case: the system checks whether pymupdf4llm's own text-layer extraction already produced a valid Markdown pipe table (detected by scanning for a \texttt{|---|---|} separator line) before deciding which pipeline to use.
If it did, the page is treated like a TEXT page: no image is rendered and no LLM call is made, since pymupdf4llm's table detector reads the native text layer directly and never paraphrases a cell value or drops the heading above the table.
Only when pymupdf4llm's extraction did not yield a usable table --- typically a scanned page or a table embedded as an image --- does the page fall back to image rendering with the vision model.
This keeps the system's total LLM cost down without sacrificing the pages that genuinely need the vision path.

Text extraction is handled by pymupdf4llm \autocite{pymupdf4llm}.
The system calls its \texttt{to\_markdown()} function once for the whole
document, with \texttt{page\_chunks=True}, which returns one result per page
as a list.
This matters for performance: pymupdf4llm runs a full document scan to identify header font sizes on every call, regardless of how many pages are requested.
Calling it once per page would repeat that scan $N$ times, making extraction quadratic in the number of pages.
Calling it once and indexing into the result list keeps extraction linear.
Any inline figures found on TEXT pages are written to a configurable output directory, and their paths are corrected to be relative to the output file.

For the hybrid strategy, a second text extraction runs in parallel using a custom font-size-based extractor built on PyMuPDF \autocite{pymupdf}.
This extractor builds a character-weighted tally of all font sizes on the page to identify the body text size, defined as the most frequent size by character count.
Each line's dominant font size is then compared to the body size as a ratio: lines at 1.8 times the body size or above are promoted to H1, 1.4 times to H2, and 1.15 times to H3.
This produces a lightweight Markdown representation of the text layer that the LLM receives alongside the rendered image.

Image rendering converts each page to a base64-encoded PNG using PyMuPDF's pixmap renderer.
The render resolution is 150 dots per inch (DPI), set via a \texttt{fitz.Matrix} scale factor of $150/72$ in both dimensions, which is about twice the PDF native resolution of 72 DPI.
The higher resolution preserves fine detail in tables, formulas, and diagrams, which improves LLM recognition accuracy.
The encoded string is passed directly in the API request body as a data URI, so no temporary files are written to disk for vision pages.
The render DPI is a configurable constant, so it can be raised for documents with very dense content or lowered to reduce memory use.
Embedded figures can also be extracted separately using \texttt{extract\_page\_figures}, which saves each image to disk and returns relative Markdown paths that the LLM can reference in its output.
The system also preserves the printed page label for each page using PyMuPDF's \texttt{get\_label()} method, which returns the label as it appears in the PDF (for example, ``iii'' for a Roman-numeral front-matter page rather than the 0-based index). These labels appear in the \texttt{<!{-}{-} Page iii -{-}>} comments that separate pages in the output file.

\section{LLM Integration}

The system uses five prompt variants, each written for a specific content type.
The \textit{default} is a general-purpose prompt used when no more specific type
is detected. The \textit{text} variant is stricter: it instructs
the model to assign heading levels by numbering depth, mapping chapter headings
to \texttt{\#}, sections to \texttt{\#\#}, and subsections to \texttt{\#\#\#}.
It exists for an optional mode in which the text path's raw output is passed
through the LLM for cleanup; in the default configuration, text pages never
reach the model at all, so this variant stays unused.
Visual pages use one of three specialised variants. The \textit{formula} variant
instructs the model to transcribe all mathematical content as \LaTeX{} notation,
preserving subscripts, superscripts, and operators exactly. The \textit{table}
variant focuses on reconstructing pipe tables, with instructions to preserve all
cell values and approximate merged cells in flat Markdown. The \textit{diagram}
variant asks the model to extract all visible text from figures and charts and
add a short description of what the diagram conveys.

The message structure sent to the API differs by strategy. Text pages go to the
API as a plain string in a single user message, with no image content. Vision
pages are sent as multipart messages: the image strategy includes one or more
base64-encoded images as \texttt{image\_url} content blocks, with no text from
the PDF. The hybrid strategy combines both, with the extracted text as the first
content block and the page image following it. All strategies pass the selected
system prompt in a system-role message before the user turn, using the
OpenAI-compatible chat completions format (see Section~\ref{sec:background:llms}).

When embedded figures have been extracted from the page in advance, as described in Section~5.3, a short text block is added to the user message listing them as Markdown image links.
This instructs the model to include those links at the appropriate positions in its output rather than describing the figures with placeholder text.
For non-English documents, a language-preservation note is appended to the system prompt by the \texttt{with\_language\_hint} function.
The note names the document language and instructs the model to keep all text in the original language.
This prevents the model from silently translating German or French content into English.

The adaptive strategy routes each page type to a specific prompt variant.
FORMULA pages use the \textit{formula} prompt, TABLE pages the \textit{table} prompt, IMAGE pages the \textit{diagram} prompt, and MIXED pages the general \textit{default} prompt.
TEXT pages are handled differently: the system calls the text strategy with \texttt{llm\_call=None}, which returns the raw pymupdf4llm output without invoking the LLM at all.
This saves the full round-trip cost on every plain-text page, which is the primary source of efficiency gain for the adaptive strategy on text-heavy documents.

Two safeguards keep this routing from losing content. Both appear in Figure~\ref{fig:routing}.
The first handles TEXT pages that are not as plain as they look.
When the pre-extracted Markdown contains an image link, the page holds an embedded raster image the classifier's image count missed, most often a formula rendered as a small picture.
Such a page is rerouted to the vision model so the formula becomes LaTeX instead of an opaque image.
The images pymupdf4llm already saved are passed along as figure references, so they are not silently dropped.
Table-of-contents pages are exempt, since the postprocessing TOC converter must receive them on the text path.

The second safeguard checks the vision model's answer before trusting it.
A response shorter than 30 per cent of the page's own extractable text almost always means the model was cut off partway, not that the page has little text.
The system then discards the vision output and falls back to the pymupdf4llm text, which is known to be complete.
This check runs after the reroute above and on every MIXED page.
A truncated model response can cost one page's formatting, but it can no longer cost the page's content.

\section{Post-processing Pipeline}
\label{sec:postprocessing_pipeline}

The raw output from pymupdf4llm and the LLM both contain predictable artefacts that must be removed before the Markdown is usable.
The pipeline applies a fixed sequence of passes in two stages.
The first stage runs per page via \texttt{clean\_page()}, before individual pages are joined into a document.
The second stage runs on the full joined document via \texttt{postprocess\_markdown()}.
Separating the two stages matters because some artefacts, such as duplicate running headers, only become visible once multiple pages are combined.

The first group removes structural noise introduced by the PDF format.
Running headers are repeated chapter or section titles that PDF places at the top of every page.
pymupdf4llm extracts these as plain-text lines; the pipeline detects them by pattern and removes them from the first line of each page chunk.

The second group removes artefacts specific to LLM output.
Decorative horizontal rules that appear directly after a heading are stripped; these come from PDF underlines misread as Markdown separators.
Prompt instruction text that some models copy verbatim into their output is removed.
\texttt{```markdown} fences that some models place around their entire response are stripped: the fence is valid syntax for a code block but not for the top-level document.

The third group targets heading quality.
pymupdf4llm often collapses all headings to the same level (\texttt{\#\#}), so the pipeline re-assigns levels based on section-number patterns: a line matching \texttt{1.1.1 Title} becomes \texttt{\#\#\#}, a line matching \texttt{1.1 Title} becomes \texttt{\#\#}, and a standalone chapter title becomes \texttt{\#}.
A separate pass merges cases where the section number and its title were emitted as two consecutive heading lines at the same level.
A further pass handles bare section numbers with no title text: the pipeline looks up the full title in the raw PyMuPDF text for the same page and patches the heading, recovering titles that the layout extraction missed.
Headings that are single words with no recognised structural keyword are demoted to plain text.

The fourth group handles special content: table-of-contents pages, list pages, figures, and encoding.
In the adaptive strategy, TOC pages are handled before per-page extraction runs.
The pipeline reads the PDF bookmark outline directly via \texttt{doc.get\_toc()} and formats it as a nested Markdown list with display page labels.
The first TOC page receives this list; continuation pages are left empty.
This bypasses pymupdf4llm entirely for TOC pages, avoiding the noisy pipe-table output it produces for these pages.
Front-matter list pages (List of Figures, List of Tables, and Listings) are classified as TEXT and processed by pymupdf4llm, then a dedicated pass converts the pipe-table structure into clean Markdown and removes dot leaders.
Figure caption lines are normalised to a single style: bold label, plain caption text.
pymupdf4llm emits captions as bold lines, and the vision model sometimes emits them fully italicised; both forms become the same bold-label style used for source-code captions.
Captions that appear before their image are swapped, blank lines between an image and its caption are removed, and label-only captions with no descriptive text are dropped.
Figure~\ref{fig:caption_reorder} shows the swap on a corpus page: the vision model placed the caption above the figure link, and the pass moves it directly below the image.
The link target is the file that \texttt{extract\_page\_figures} saved to disk (Section~5.3), so this is also what a preserved figure looks like in the final Markdown.
Italic spans containing only Greek letters are converted to inline LaTeX (e.g.\ an italic $\varphi$ becomes \texttt{\$\textbackslash varphi\$}).
OCR superscript artefacts of the form \texttt{\_EMC\_[2]} are replaced with Unicode superscript characters.

\begin{figure}[ht]
\medskip\noindent\textbf{Raw vision-model output}
\begin{lstlisting}[language={},nolol,identifierstyle=\color{black},basicstyle=\footnotesize\ttfamily\color{black},upquote=true,breakatwhitespace=true]
**Figure 1.1:** What is this Thesis about?

![Figure 1](figures/page_009_fig_001.png)
\end{lstlisting}

\medskip\noindent\textbf{After caption reordering}
\begin{lstlisting}[language={},nolol,identifierstyle=\color{black},basicstyle=\footnotesize\ttfamily\color{black},upquote=true,breakatwhitespace=true]
![Figure 1](figures/page_009_fig_001.png)
**Figure 1.1:** What is this Thesis about?
\end{lstlisting}
\caption{Caption reordering on a corpus page (sentiment analysis thesis,
page 9, classified as MIXED). The vision model emits the caption above the
figure link. The pass swaps the two lines and removes the blank line, so
the caption sits directly below its image. The link path is where
\texttt{extract\_page\_figures} saved the embedded figure.}
\label{fig:caption_reorder}
\end{figure}

Source-code listings needed a detection path of their own.
\LaTeX's \texttt{listings} package sets code in a fixed-width font.
pymupdf4llm does not read that as a signal to build a code block.
It either leaves the listing as plain text, or, when a numbered line looks like a heading, breaks it into fragments wrapped in scattered backticks.

The pipeline finds monospace lines directly from PyMuPDF's font data.
A font counts as monospace by name, matching substrings such as Courier or Consolas, or by measurement: if every character on the page has the same width, the font is fixed-width no matter what it is called.
This second check exists because some PDFs give their code font a generic internal name with no clue in it at all.
A page's lines are only trusted as code if at least one shows real code syntax, such as a line-number gutter or a closing parenthesis before a colon.
This rule keeps out bibliography entries and citation URLs, which are sometimes set in the same font as a listing but are not code.

A single page can hold more than one listing with prose between them.
The reconstruction scans the page line by line and only groups consecutive matching lines into one fence.
Prose and unrelated listings in between stay untouched.
This stops two listings from merging into each other or swallowing the caption between them.
Line numbers are stripped and renumbered in order, since \LaTeX{} only prints a number for a line's first physical row.
A wrapped continuation gets no number at all, which looks broken once reflowed into one fence.
Figure~\ref{fig:code_reconstruction} shows the pass on a corpus page.
pymupdf4llm merges pairs of code lines into list items and leaves the line numbers inline.
The reconstruction returns the listing as a fenced block, one code line per line, and the caption below it stays outside the fence.

\begin{figure}[ht]
\medskip\noindent\textbf{Raw pymupdf4llm output}
\begin{lstlisting}[language={},nolol,identifierstyle=\color{black},basicstyle=\footnotesize\ttfamily\color{black},literate={ü}{{\"u}}1 {ä}{{\"a}}1,upquote=true,breakatwhitespace=true]
- 1 def getAction(state): # gibt -1 für Schläger links oder +1 für rechts zurück 2 global epsilon, Q

- 3 if np.random.rand() <= epsilon: 4 return np.random.choice([-1, 1])

- 5 return (np.argmax(Q[int(state)]) * 2) - 1

**Quellcode 5.2:** Methode getAction [8, S. 283]
\end{lstlisting}

\medskip\noindent\textbf{After code block reconstruction}
\begin{lstlisting}[language={},nolol,identifierstyle=\color{black},basicstyle=\footnotesize\ttfamily\color{black},literate={ü}{{\"u}}1 {ä}{{\"a}}1,upquote=true,breakatwhitespace=true]
```
1 def getAction(state): # gibt -1 für Schläger links oder +1 für rechts zurück
2 global epsilon, Q
3 if np.random.rand() <= epsilon:
4 return np.random.choice([-1, 1])
5 return (np.argmax(Q[int(state)]) * 2) - 1
```

**Quellcode 5.2:** Methode getAction [8, S. 283]
\end{lstlisting}
\caption{Source-code reconstruction on a corpus page (neural networks
thesis, page 21). pymupdf4llm emits the listing as list items that merge
two code lines each; the monospace-font pass rebuilds it as a fenced code
block and leaves the caption line untouched.}
\label{fig:code_reconstruction}
\end{figure}

The same monospace check is reused during figure extraction.
\LaTeX's \texttt{listings} package draws a shaded box behind a listing.
That box is a vector-filled shape, not a raster image.
Earlier, \texttt{extract\_page\_figures} treated any large vector-drawn region as a possible figure.
This mistook the shaded box for a diagram and saved a screenshot of the code instead of leaving it as text.
The function now checks whether a candidate region's text is mostly monospace before treating it as a figure, and skips it if so.

A second problem appears when a page holds more than one vector-drawn element, such as two separate code boxes.
Taking the bounding box of every vector drawing on the page and merging them into one region spans all the prose in between.
\texttt{extract\_page\_figures} now groups vector drawings by vertical distance and checks each group on its own.
Unrelated regions on the same page no longer merge into one oversized crop.

A last check filters out pymupdf4llm's own fallback for complex equations.
When it cannot read a formula as text, it renders that formula as a small image instead, usually under 50 pixels tall for one line of math.
Carrying that image over to the vision model would duplicate the formula, since the model already converts it to \LaTeX{} correctly from the page image.
The height check drops the image before that duplication happens.

The second stage handles artefacts that only become visible on the joined document.
Duplicate running headers that cross page boundaries are removed a second time.
Lone page-number paragraphs, whether Arabic digits or Roman numerals, are stripped wherever they appear.
Title-page metadata headings are demoted, and chapter headings split across pages are merged.
Two further demotion passes run here: headings whose entire content is italic, and headings that are really numbered source-code listing lines, both become plain or list text.
OCR text blocks that pymupdf4llm extracts from embedded figures are removed, since the figure is already included as a Markdown image link and the OCR output adds only noise.
LaTeX math delimiters in the \texttt{\textbackslash(...\textbackslash)} and \texttt{\textbackslash[...\textbackslash]} styles are normalised to the \texttt{\$...\$} and \texttt{\$\$...\$\$} forms used throughout the output.

List pages get one more pass at this stage.
pymupdf4llm sometimes renders the same list of figures or tables as plain bullets on one page and as table rows on the next, since it lays out each page on its own.
A page-crossing pass rewrites the bullet entries into the same row format, so the whole list reads the same way from start to end.
A second pass then merges table rows that a wrapped entry split into several pieces, joining them back into one row per entry and reversing any hyphenation that broke mid-word at the wrap point.

Figure~\ref{fig:postprocessing} shows all passes in both stages.

\begin{figure}[H]
    \centering
    % Phase 1: classification cascade — mirrors analyze_page() in strategies/adaptive.py
\begin{tikzpicture}[
    font=\sffamily\small,
    node distance=4.5mm and 9mm,
    >={Latex[length=1.8mm]},
    box/.style={draw=black!70, line width=0.5pt, rounded corners=2pt, align=center, inner sep=5pt},
    signal/.style={box, fill=black!8, text width=78mm},
    cond/.style={box, fill=white, font=\ttfamily\footnotesize, align=left, text width=78mm},
    outcome/.style={box, fill=black!12, font=\sffamily\footnotesize\bfseries, text width=27mm, rounded corners=5pt},
    lbl/.style={font=\sffamily\scriptsize\itshape, inner sep=2pt}
]
\node[signal] (sig) {extract signals per page:\\[2pt]
    {\ttfamily\footnotesize image\_count \(\cdot\) vector\_path\_count\\[1pt] text\_length \(\cdot\) table\_count}};

\node[cond, below=of sig] (c1) {text\_length < 10 and not has\_images\\and vector\_path\_count < 5};
\node[cond, below=of c1]  (c2) {has\_formulas and text\_length < 50};
\node[cond, below=of c2]  (c3) {has\_images and image\_count > 3};
\node[cond, below=of c3]  (c4) {table\_count > 0};
\node[cond, below=of c4]  (c5) {has\_images or has\_formulas};
\node[cond, below=of c5]  (c6) {text\_length >= 50};

\node[outcome, right=of c1] (o1) {EMPTY \(\cdot\) 0.90};
\node[outcome, right=of c2] (o2) {FORMULA \(\cdot\) 0.85};
\node[outcome, right=of c3] (o3) {IMAGE \(\cdot\) 0.90};
\node[outcome, right=of c4] (o4) {TABLE \(\cdot\) 0.85};
\node[outcome, right=of c5] (o5) {MIXED \(\cdot\) 0.75};
\node[outcome, right=of c6] (o6) {TEXT \(\cdot\) 0.95};
\node[outcome, below=of c6] (o7) {MIXED \(\cdot\) 0.50};

\draw[->] (sig) -- (c1);
\draw[->] (c1) -- node[lbl, above] {yes} (o1);
\draw[->] (c2) -- node[lbl, above] {yes} (o2);
\draw[->] (c3) -- node[lbl, above] {yes} (o3);
\draw[->] (c4) -- node[lbl, above] {yes} (o4);
\draw[->] (c5) -- node[lbl, above] {yes} (o5);
\draw[->] (c6) -- node[lbl, above] {yes} (o6);
\draw[->] (c1) -- node[lbl, right] {no} (c2);
\draw[->] (c2) -- node[lbl, right] {no} (c3);
\draw[->] (c3) -- node[lbl, right] {no} (c4);
\draw[->] (c4) -- node[lbl, right] {no} (c5);
\draw[->] (c5) -- node[lbl, right] {no} (c6);
\draw[->] (c6) -- node[lbl, right] {no} (o7);
\end{tikzpicture}

    \caption[Page classifier, phase 1: the classification cascade]{Page classifier, phase 1: the classification cascade in \texttt{analyze\_page}. Conditions are checked top to bottom; the first match assigns the type and its confidence score.}
    \label{fig:classifier}
\end{figure}

\begin{figure}[H]
    \centering
    % Phase 2: routing checks — mirrors adaptive_strategy() in strategies/adaptive.py
\begin{tikzpicture}[
    font=\sffamily\small,
    >={Latex[length=1.8mm]},
    box/.style={draw=black!70, line width=0.5pt, rounded corners=2pt, align=center, inner sep=4pt},
    head/.style={box, fill=black!12, font=\sffamily\footnotesize\bfseries, rounded corners=5pt},
    cond/.style={box, fill=white, text width=38mm, minimum height=14mm, font=\sffamily\footnotesize},
    outcome/.style={box, fill=black!5, text width=21mm, font=\sffamily\footnotesize, anchor=north},
    direct/.style={align=center, font=\sffamily\footnotesize},
    lbl/.style={font=\sffamily\scriptsize\itshape, inner sep=2pt}
]
% --- column 1: TABLE ---
\node[head] (h1) at (0,0) {TABLE \(\cdot\) 0.85};
\node[cond] (q1) at (0,-2.1) {text layer already gives a valid pipe table?};
\node[outcome] (a1) at (-1.25,-3.7) {text path:\\no LLM call};
\node[outcome] (b1) at (1.25,-3.7) {vision LLM,\\table prompt};
\draw[->] (h1) -- (q1);
\draw[->] (q1.south) -- node[lbl, left=1pt] {yes} (a1.north);
\draw[->] (q1.south) -- node[lbl, right=1pt] {no} (b1.north);

% --- column 2: TEXT ---
\node[head] (h2) at (5.5,0) {TEXT \(\cdot\) 0.95};
\node[cond] (q2) at (5.5,-2.1) {extracted Markdown contains image links?\\{\scriptsize(TOC pages exempt)}};
\node[outcome] (a2) at (4.25,-3.7) {reroute:\\MIXED \(\cdot\) 0.75};
\node[outcome] (b2) at (6.75,-3.7) {text path:\\no LLM call};
\draw[->] (h2) -- (q2);
\draw[->] (q2.south) -- node[lbl, left=1pt] {yes} (a2.north);
\draw[->] (q2.south) -- node[lbl, right=1pt] {no} (b2.north);

% --- column 3: MIXED ---
\node[head] (h3) at (11,0) {MIXED \(\cdot\) any score};
\node[cond] (q3) at (11,-2.1) {vision response shorter than 30\,\% of extracted text?};
\node[outcome] (a3) at (9.75,-3.7) {fallback:\\pymupdf4llm text};
\node[outcome] (b3) at (12.25,-3.7) {keep vision output};
\draw[->] (h3) -- (q3);
\draw[->] (q3.south) -- node[lbl, left=1pt] {yes} (a3.north);
\draw[->] (q3.south) -- node[lbl, right=1pt] {no} (b3.north);

% --- direct routes (no check) ---
\draw[black!50, dashed, thin] (-2.4,-5.4) -- (13.4,-5.4);
\node[direct, font=\sffamily\footnotesize\itshape, anchor=north] at (5.5,-5.65) {direct routes, no check:};
\node[direct, anchor=north] at (0,-6.35)   {EMPTY \(\rightarrow\) skipped};
\node[direct, anchor=north] at (5.5,-6.35) {FORMULA \(\rightarrow\) vision LLM,\\formula prompt};
\node[direct, anchor=north] at (11,-6.35)  {IMAGE \(\rightarrow\) vision LLM,\\diagram prompt};
\end{tikzpicture}

    \caption[Page classifier, phase 2: the routing checks]{Page classifier, phase 2: the routing checks in \texttt{adaptive\_strategy}. TABLE pages skip the LLM when the text layer already holds a valid pipe table. TEXT pages with image links are rerouted to MIXED. Every MIXED page gets a length check on the vision response. EMPTY, FORMULA, and IMAGE pages route directly without a check.}
    \label{fig:routing}
\end{figure}

\begin{figure}[H]
    \centering
    % Two-stage post-processing — mirrors clean_page() and postprocess_markdown() in postprocess.py
\begin{tikzpicture}[
    font=\sffamily\small,
    >={Latex[length=1.8mm]},
    box/.style={draw=black!70, line width=0.5pt, rounded corners=2pt, align=center, inner sep=4pt},
    startstop/.style={box, fill=black!12, rounded corners=5pt, font=\sffamily\footnotesize\bfseries},
    pass/.style={box, fill=white, font=\sffamily\footnotesize, text width=48mm},
    stagetitle/.style={align=center, font=\sffamily\footnotesize},
    stageframe/.style={draw=black!50, dashed, thin, rounded corners=3pt, inner sep=3mm}
]
% --- Stage 1: clean_page, per page (left column) ---
\node[stagetitle] (t1) at (0,0) {\textbf{Stage 1 \(\cdot\) per page}\\\texttt{clean\_page}};
\node[pass, below=3.5mm of t1] (b1) {structural noise\\{\scriptsize running headers \(\cdot\) decorative\\rules under headings}};
\node[pass, below=3.5mm of b1] (b2) {LLM artefacts\\{\scriptsize \texttt{```markdown} fences \(\cdot\)\\echoed prompt text}};
\node[pass, below=3.5mm of b2] (b3) {heading quality\\{\scriptsize level normalisation \(\cdot\) split-heading\\merge \(\cdot\) bare-number recovery \(\cdot\)\\demotions}};
\node[pass, below=3.5mm of b3] (b4) {special content\\{\scriptsize TOC tables \& dot leaders \(\cdot\)\\list-of-figures pages \(\cdot\) figure captions\\Greek \(\rightarrow\) LaTeX \(\cdot\) OCR superscripts}};
\node[pass, below=3.5mm of b4] (b5) {source-code listings\\{\scriptsize fences rebuilt from monospace\\lines \(\cdot\) renumbering}};
\node[stageframe, fit=(t1)(b1)(b2)(b3)(b4)(b5)] (S1) {};

% --- Stage 2: postprocess_markdown, full document (right column) ---
\node[stagetitle] (t2) at (8.8,0) {\textbf{Stage 2 \(\cdot\) full document}\\\texttt{postprocess\_markdown}};
\node[pass, below=3.5mm of t2] (d1) {cross-page duplicates\\{\scriptsize running headers \(\cdot\) page numbers \(\cdot\)\\repeated headings}};
\node[pass, below=3.5mm of d1] (d2) {heading demotions \& merges\\{\scriptsize title page \(\cdot\) italic \& listing\\headings \(\cdot\) chapter merge}};
\node[pass, below=3.5mm of d2] (d3) {content cleanup\\{\scriptsize OCR picture text \(\cdot\) bibliography\\entries \(\cdot\) LaTeX delimiter\\normalisation}};
\node[pass, below=3.5mm of d3] (d4) {list-page consistency\\{\scriptsize bullets \(\rightarrow\) table rows \(\cdot\)\\wrapped-row merge}};
\node[stageframe, fit=(t2)(d1)(d2)(d3)(d4)] (S2) {};

% --- endpoints and join ---
\node[startstop, above=4mm of S1] (raw) {raw page Markdown};
\node[box, fill=black!5, font=\sffamily\footnotesize, text width=22mm, align=center] (join) at (4.4,-4.4) {join all pages\\with \texttt{-{}-{}-}\\separator};
\node[startstop, below=4mm of S2] (final) {final Markdown document};

% --- flow ---
\draw[->] (raw) -- (S1);
\draw[->] (S1.east|-join) -- (join.west);
\draw[->] (join.east) -- (S2.west|-join);
\draw[->] (S2) -- (final);
\end{tikzpicture}

    \caption[Two-stage post-processing pipeline]{Two-stage post-processing pipeline: \texttt{clean\_page} applies structural, LLM artefact, heading, and special-content passes per page; \texttt{postprocess\_markdown} applies heading, running-header, and encoding passes to the full joined document}
    \label{fig:postprocessing}
\end{figure}

\section{Implementation Details}

The system is implemented in Python and depends on four external libraries.
PyMuPDF \autocite{pymupdf} (version 1.27) provides the core PDF interface: page analysis, text extraction, image rendering, and table detection.
pymupdf4llm \autocite{pymupdf4llm} (version 1.27, matching PyMuPDF) wraps it to produce Markdown output directly from the text layer.
The \texttt{requests} library handles all HTTP communication with LM Studio.
\texttt{langdetect} detects the document language once from the first page's extracted text. Since thesis documents use a consistent language throughout, the same hint is applied to all LLM calls in that document.
No training or fine-tuning is required. The system runs entirely with pre-trained models loaded through LM Studio \autocite{lmstudio}.

A single model is used across all LLM-dependent strategies, as described in
Section~4.4: Qwen2.5-VL~7B, a vision-language model served through LM Studio,
which exposes an OpenAI-compatible \texttt{/v1/chat/completions} endpoint. The
text strategy and the adaptive TEXT path do not invoke the LLM at all. The
server URL and model name are passed as parameters to every strategy
function, so switching to a different model, or to a server running on a
remote machine instead of the same host, requires no code changes.

Reproducibility is ensured through three design choices. All experiments run at
temperature 0.0, which makes LLM output deterministic for the same input.
Library versions are pinned in the requirements file, so extraction and rendering
behaviour is identical across runs. Every LLM call goes through the injectable
\texttt{call\_llm} parameter, which means the full pipeline can be exercised in
tests using a fixed-response stub, without any dependency on a running LM Studio
instance.
The evaluation framework stores all results as JSON and generates reports from that file, so results can be re-analysed without re-running the experiment.
